% =========================
% Cas pratique 3
% =========================

\section[Cas 3]{Cas pratique 3 : Confidentialité en
contrefaçon}

%--------------------------
\begin{frame}{Question 1: hypothèse rien de prévu}

La question propose le découpage suivant: 
Ensemble de documents "informations" = documents pas confidentiels $\cup$
documents confidentiels.\\
Objet: documents pas confidentiels.

\begin{block}{}
Interprétation stricte de l'ordonnance (\textbf{pourvoi n°03-12162}):
appartenait à l'huissier d'effectuer les opérations de saisie 
conformément aux dispositions de l'ordonnance le saisissant
\end{block}

\begin{itemize}
    \item En l'espèce,
    \begin{itemize}
        \item Hypothèse 1: saisie documents pas confidentiels $\in$ ordonnance.
        Non, pas de mention dans le sujet.
        \item Hypothèse 2: saisie documents pas confidentiels $\notin$ ordonnance.
        Oui, car rien de prévu (d'autre) dans l'ordonnance.
        Or, documents pas confidentiels $\neq$ échantillons de tube.
    \end{itemize}
    \item En conclusion, saisie illicite. 
    La saisie est licite si l'ordonnance vérifie l'hypothèse 1.
\end{itemize}

\end{frame}
%--------------------------

%--------------------------
\begin{frame}{Question 1: hypothèse séquestre provisoire}

Objet: documents pas confidentiels.

\begin{block}{}
Art. \textbf{R615-4} al. 4 CPI dispose que "Afin d'assurer la protection du secret des affaires, 
le président peut ordonner d'office le placement sous séquestre provisoire des pièces saisies, 
dans les conditions prévues à l'article R. 153-1 du code de commerce.".
\end{block}

\begin{itemize}
    \item En l'espèce, documents pas confidentiels.
    \begin{itemize}
        \item Hypothèse 1: saisie document pas confidentiels $\in$ séquestre provisoire. 
        \item Hypothèse 2: saisie document pas confidentiels $\notin$ séquestre provisoire.
    \end{itemize}
    Aucun élément ne permet de discrimer entre les deux hypothèses.
    \item En conclusion,
    \begin{itemize}
        \item Conclusion 1: la saisie est illicite.
        \item Conclusion 2: la saisie est licite.
    \end{itemize}
    La saisie est licite si le séquestre provisoire vérifie l'hypothèse 2.
\end{itemize}

\end{frame}
%--------------------------

%--------------------------
\begin{frame}{Question 2: hypothèse rien de prévu}

Ensemble de documents = documents pertinents $\cup$ documents pas pertinents.\\
Objet: ensemble de documents.

\begin{block}{}
Art. \textbf{L615-5} al. 2 3e phrase CPI dispose que "L'ordonnance peut autoriser 
[ci-après: première condition (i)] la saisie réelle de tout document se rapportant aux produits ou procédés prétendus contrefaisants
[ci-après: deuxième condition (ii)] en l'absence de ces derniers.".
\end{block}

\begin{itemize}
    \item On suppose que l'ordonannce a prévu ce cas-là (sinon 
    saisie illicite). Le stock de tube
        contrefaisant est 50\% de son chiffre d'affaire.
        Donc, ensemble de documents comprend probablement
        des documents pas pertinents. Donc,
        documents pas pertinents $\in$ saisie.
        Donc, première condition (i) pas vérifiée.
        De plus, il y a des stocks de tubes.
        Donc, deuxième condition (ii) pas vérifiée.

    \item En conclusion, la saisie est illicite.
    Pour s'oppposer à la saisie pendant la saisie-cf., 
    vérifier que l'ordonnance a prévu ce cas-là.
    Pour contester ensuite auprès du juge,
    vérifier la première condition (i)
    et vérifier la deuxième condition (ii).
\end{itemize}

\end{frame}

%--------------------------

%--------------------------

\begin{frame}{Question 2: hypothèse séquestre provisoire}


Objet: ensemble de documents.

\begin{block}{}
Art. \textbf{R615-4} al. 4 CPI dispose que "Afin d'assurer la protection du secret des affaires, 
le président peut ordonner d'office le placement sous séquestre provisoire des pièces saisies, 
dans les conditions prévues à l'article R. 153-1 du code de commerce.".
\end{block}

\begin{itemize}
    \item En l'espèce,
    ensemble de documents = documents pas confidentiels $\cup$ documents confidentiels.
    \item En conclusion, pas d'opposition pendant la saisie-cf.
    Toutefois, si documents pas confidentiels $\in$ séquestre provisoire,
    alors, l'ordonnance est nulle.\\    
    Donc, on comprend que le breveté n'a aucun intérêt à demander
    un séquestre provisoire dans l'ordonnance.
\end{itemize}

\end{frame}
%--------------------------

%--------------------------
\begin{frame}{Question 3}

\begin{block}{}
Art. \textbf{R153-1} al. 2 Code de Commerce: Si le juge n'est pas saisi d'une demande
 de modification ou de rétractation 
de son ordonnance en application de l'article 497 du code de procédure civile 
dans un délai d'un mois à compter de la signification de la décision, 
la mesure de séquestre provisoire mentionnée à l'alinéa précédent est levée 
et les pièces sont transmises au requérant.
\end{block}

\begin{itemize}
    \item Hypothèse rien de prévu: Asie Trading notifie le commissaire
    de justice que les documents sont confidentiels. Les documents
    sont placé au séquestre provisoire.
    \item Hypothèse séquestre provisoire: Les documents
    sont placés d'office au séquestre provisoire.
\end{itemize}
Dans les deux cas, la partie saisie doit saisir le juge dans un délai de 1 mois.
    Sinon, il y a levée du sequestre provisoire.
\end{frame}

%--------------------------